\section{Intégration numérique : le retour radial}\label{sec:integration}

\subsection{Le problème incrémental}

Le code EF découpe le chargement en pas $[t_n,t_{n+1}]$. En un point de Gauss, on connaît l'état au début du pas $(\epsp_n, p_n, \X_n)$ et la déformation totale en fin de pas $\eps_{n+1}$, que propose l'itération courante de l'équilibre global. On cherche $(\sig_{n+1},\epsp_{n+1},p_{n+1},\X_{n+1})$. C'est l'interface implémentée :
\begin{lstlisting}
sigma_new, state_new = integrate(eps_new, state_old, params)   # state = (eps_p, p, alpha)
\end{lstlisting}

\subsection{Discrétisation par Euler implicite}

On intègre les lois d'évolution~\eqref{eq:flow} par le schéma d'Euler \emph{implicite} (rétrograde), qui évalue la direction d'écoulement en fin de pas :
\begin{align}
\epsp_{n+1} &= \epsp_n + \Delta p\,\N_{n+1}, & \X_{n+1} &= \X_n + \tfrac23 C\,\Delta p\,\N_{n+1},\\
p_{n+1} &= p_n+\Delta p, & f_{n+1} &\le 0,\quad \Delta p\ge0,\quad \Delta p\,f_{n+1}=0.
\end{align}
Le schéma implicite~\cite{simo1998,desouza2008} est inconditionnellement stable et garantit que l'état final est admissible ($f_{n+1}\le0$). L'inconnue principale est le scalaire $\Delta p\ge0$.

\subsection{Prédicteur élastique}

On suppose d'abord le pas élastique ($\Delta p=0$) :
\begin{equation}
\sig\tri = \C:(\eps_{n+1}-\epsp_n),\qquad \s\tri = 2\mu\dev(\eps_{n+1}-\epsp_n),\qquad \bs\xi\tri = \s\tri-\X_n,\qquad q\tri = \sqrt{\tfrac32\,\bs\xi\tri:\bs\xi\tri}.
\end{equation}
Si $f\tri = q\tri - R(p_n)\le0$, l'hypothèse est vérifiée et le pas est élastique. Sinon il faut corriger.

\subsection{Correcteur plastique}

\begin{proposition}[Colinéarité]
En cas de plastification, $\bs\xi_{n+1}$ est colinéaire à $\bs\xi\tri$, et $q_{n+1} = q\tri - (3\mu+C)\Delta p$.
\end{proposition}
\begin{proof}
On a $\s_{n+1} = \s\tri - 2\mu\Delta p\N_{n+1}$ (la correction n'agit que sur le déviateur, car $\N$ est déviatorique), et $\X_{n+1} = \X_n + \frac23C\Delta p\N_{n+1}$. Par différence :
\[
\bs\xi_{n+1} = \bs\xi\tri - \big(2\mu+\tfrac23C\big)\Delta p\,\N_{n+1} = \bs\xi\tri - \big(2\mu+\tfrac23C\big)\Delta p\,\frac32\frac{\bs\xi_{n+1}}{q_{n+1}}.
\]
Donc $\bs\xi_{n+1}\big(1+\frac{(3\mu+C)\Delta p}{q_{n+1}}\big) = \bs\xi\tri$. Le facteur est un scalaire positif, d'où la colinéarité. En prenant $\sigma_{eq}$ des deux membres : $q_{n+1}+(3\mu+C)\Delta p = q\tri$.
\end{proof}

Conséquence importante : $\N_{n+1} = \N\tri = \frac32\bs\xi\tri/q\tri$ est connu dès le prédicteur. Le retour se fait le long du « rayon » qui joint l'état d'essai au centre du domaine, d'où le nom de \emph{retour radial}. Le problème tensoriel se réduit à \textbf{une équation scalaire}, la condition $f_{n+1}=q_{n+1}-R(p_{n+1})=0$ :
\begin{equation}
\boxed{\begin{aligned}
g(\Delta p) \;=\; q\tri - (3\mu+C)\,\Delta p - R(p_n+\Delta p) \;=\; 0
\end{aligned}}
\label{eq:g}
\end{equation}

\paragraph{Écrouissage linéaire.} $R$ est affine et~\eqref{eq:g} donne directement
\[
\Delta p = \frac{f\tri}{3\mu+C+H}.
\]
Pour $C=0$, on retrouve la formule du sujet, $\Delta p = f\tri/(3\mu+H)$.

\paragraph{Écrouissage de Voce.} $g$ est non linéaire. On la résout par Newton :
\[
\Delta p^{(i+1)} = \Delta p^{(i)} - \frac{g(\Delta p^{(i)})}{g'(\Delta p^{(i)})},\qquad g'(\Delta p) = -\big(3\mu+C+R'(p_n+\Delta p)\big)<0.
\]
On l'initialise avec $\Delta p^{(0)}=f\tri/(3\mu+C+R'(p_n))$, qui est exact dans le cas linéaire. Comme $R$ est concave ($R''=-Qb^2e^{-bp}\le0$), $g$ est convexe et décroissante. Newton converge alors de façon monotone vers l'unique racine positive. En pratique 2 itérations suffisent pour atteindre l'arrondi machine (figure~\ref{fig:tangent}, à droite). Le code en fait un nombre \emph{fixe} (6 par défaut), parce que \code{jit} et \code{vmap} exigent un flot de contrôle statique (section~\ref{sec:jax}).

\subsection{Mise à jour}

Avec $\N = \frac32\bs\xi\tri/q\tri$ :
\begin{equation}
\sig_{n+1} = \sig\tri - 2\mu\,\Delta p\,\N,\quad \epsp_{n+1}=\epsp_n+\Delta p\,\N,\quad p_{n+1}=p_n+\Delta p,\quad \X_{n+1}=\X_n+\tfrac23C\,\Delta p\,\N.
\end{equation}
Pour $C=0$ on a $\s_{n+1}=\big(1-\frac{3\mu\Delta p}{q\tri}\big)\s\tri$, soit exactement la formule du sujet.

\begin{algorithm}[h]
\caption{Retour radial J2 (fonction \code{integrate}, sans branchement Python)}\label{alg:rr}
\begin{algorithmic}[1]
\Require $\eps_{n+1}$, $(\epsp_n,p_n,\X_n)$, paramètres
\State $\eps^e\gets\eps_{n+1}-\epsp_n$;\quad $\s\tri\gets2\mu\dev\eps^e$;\quad $\sig\tri\gets K\tr\eps^e\,\I+\s\tri$;\quad $\bs\xi\gets\s\tri-\X_n$
\State $\chi\gets\bs\xi\cdot\bs\xi$;\quad $\chi_s\gets\mathtt{where}(\chi>0,\chi,1)$;\quad $q\tri\gets\mathtt{where}(\chi>0,\sqrt{1.5\chi_s},0)$ \Comment{double where}
\State $\mathrm{plast}\gets q\tri-R(p_n)>0$
\State $q_s\gets\mathtt{where}(\mathrm{plast},q\tri,R(p_n))$ \Comment{branche élastique : solution $\Delta p=0$}
\State $\Delta p\gets$ Newton sur $g(\Delta p;q_s,p_n)=0$, nombre d'itérations fixe \Comment{dérivée implicite}
\State $\Delta p\gets\mathtt{where}(\mathrm{plast},\Delta p,0)$;\quad $\N\gets1.5\,\bs\xi/\mathtt{where}(\chi>0,q\tri,1)$
\State \Return $\sig\tri-2\mu\Delta p\N$,\ $(\epsp_n+\Delta p\N,\ p_n+\Delta p,\ \X_n+\frac23C\Delta p\N)$
\end{algorithmic}
\end{algorithm}
